% Editorial master: guides/pronunciation/answer-key.md

% Master SHA-256: 1f1f1191fb20e0e16b2c6d9f1d6d98311254f071eb86b20bdf6ce53ec3a65616

% Initial assembly; retain complete reviewed TeX for reproduction.

\newcommand{\runningtitle}{Pronunciation key}

% Layout inherited from the reviewed previous course; set a running title per unit.
% This is a small hand-maintained TeX layout, not a conversion pipeline.
\RequirePackage{fix-cm}
\documentclass[11pt,twoside]{article}
\usepackage{fontspec}
\setmainfont{Cambria}[
  Renderer=HarfBuzz,
  BoldFont=cambriab.ttf,
  ItalicFont=cambriai.ttf,
  BoldItalicFont=cambriaz.ttf
]
\setmonofont{lmmono10-regular.otf}[BoldFont=lmmonolt10-bold.otf]
\usepackage[paperwidth=165mm,paperheight=241mm,inner=21mm,outer=19mm,
  top=22mm,bottom=20mm,headheight=14pt,headsep=7mm,footskip=10mm]{geometry}
\usepackage{xcolor,array,booktabs,longtable,tabularx}
\usepackage{enumitem,fancyhdr}
\usepackage[unicode,hidelinks]{hyperref}
\definecolor{ink}{HTML}{000000}
\definecolor{muted}{HTML}{4C5558}
\setlength{\parindent}{0pt}
\setlength{\parskip}{3.5pt}
\AtBeginDocument{\fontsize{11}{14.3}\selectfont}
\setlength{\tabcolsep}{4pt}
\renewcommand{\arraystretch}{1.10}
\setlist[itemize]{leftmargin=1.3em,label=\textbullet,itemsep=2pt,topsep=3pt,parsep=0pt}
\setlist[enumerate]{leftmargin=1.8em,itemsep=3pt,topsep=4pt,parsep=0pt}
\widowpenalty=10000
\clubpenalty=10000
\emergencystretch=1em
\raggedbottom
\pagestyle{fancy}
\fancyhf{}
\providecommand{\runningtitle}{Old English}
\fancyhead[LE,RO]{\small\itshape \runningtitle}
\fancyhead[LO,RE]{\small Old English}
\fancyfoot[LE,RO]{\small\thepage}
\renewcommand{\headrulewidth}{0.3pt}
\newcommand{\booktitle}[2]{
  {\fontsize{9}{11}\selectfont\color{muted}#1\par}
  \vspace{3pt}
  {\fontsize{20}{24}\selectfont\color{ink}#2\par}
  \vspace{6pt}
}
\newcommand{\source}[1]{\par{\footnotesize\color{muted}#1\par}}
\newcommand{\readingtext}[1]{
  \par\addvspace{5pt}
  \begingroup\leftskip=1.5em
    \fontsize{12}{16}\selectfont #1\par
  \endgroup\addvspace{5pt}
}
\makeatletter
\renewcommand{\section}{\@startsection{section}{1}{\z@}
  {-2.4ex \@plus -0.6ex \@minus -0.2ex}{0.8ex \@plus 0.2ex}
  {\normalfont\fontsize{12}{15}\selectfont\bfseries\color{ink}}}
\renewcommand{\subsection}{\@startsection{subsection}{2}{\z@}
  {-1.8ex \@plus -0.5ex \@minus -0.2ex}{0.5ex \@plus 0.2ex}
  {\normalfont\fontsize{11}{14}\selectfont\bfseries\color{ink}}}
\makeatother


\hyphenpenalty=10000

\exhyphenpenalty=10000

\emergencystretch=3em

\newcommand{\reserve}[1]{\par\penalty-100\begingroup\dimen0=\pagegoal\advance\dimen0 by-\pagetotal\ifdim\dimen0<#1\newpage\fi\endgroup}

\begin{document}

\addfontfeatures{Ligatures=TeXOff}

\raggedright

\booktitle{SEPARATE KEY}{Pronunciation practice: key}

% mdparagraph:P001:begin

These answers follow the guide's chosen reading convention. The English words identify particular sounds; they are not replacement spellings for whole Old English words. The permitted American \textbf{r} and practical short-vowel approximations are acceptable. Different scholarly reconstructions do not make all sound choices interchangeable within these exercises.

% mdparagraph:P001:end



\reserve{5\baselineskip}

% mdparagraph:P002:begin

\section*{A. Length contrasts}

% mdparagraph:P002:end



% answer: PRON-P01a

% mdparagraph:P003:begin

\textbf{A1. ac, “but”; āc, “oak”.} \textbf{Ā} in \textbf{āc} is long. Both vowels have the \emph{father} quality; hold the second longer. Both final consonants are \textbf{k}. Do not replace \textbf{ā} with the English letter name \textbf{a}.

% mdparagraph:P003:end



% answer: PRON-P01b

% mdparagraph:P004:begin

\textbf{A2. is, “is”; īs, “ice”.} \textbf{Ī} in \textbf{īs} is long. Use \emph{bit} for short \textbf{i}, held \emph{see} for long \textbf{ī}. The selected practical vowels differ in quality as well as length. Both words end in unvoiced \textbf{s}, as in \emph{sip}; modern English \emph{is} does not determine Old English voicing.

% mdparagraph:P004:end



% answer: PRON-P01c

% mdparagraph:P005:begin

\textbf{A3. god, “god”; gōd, “good”.} \textbf{Ō} in \textbf{gōd} is long. The vowel of \textbf{god} is short and rounded, as described in the guide; \textbf{gōd} has the steady initial vowel of \emph{go}, held without the glide. Retain hard \textbf{g} in both. Do not copy the vowels of the modern English glosses.

% mdparagraph:P005:end



% answer: PRON-P01d

% mdparagraph:P006:begin

\textbf{A4. full, “full”; fūl, “vile”.} \textbf{Ū} in \textbf{fūl} is long, with the \emph{food} quality; \textbf{full} has short \emph{foot}-type \textbf{u}. The target is this vowel-length and vowel-quality contrast. Both final \textbf{l} sounds have ordinary single-consonant duration; the written \textbf{ll} of \textbf{full} does not require extra holding.

% mdparagraph:P006:end



\reserve{5\baselineskip}

% mdparagraph:P007:begin

\section*{B. Consonants and unfamiliar sounds}

% mdparagraph:P007:end



% answer: PRON-P02a

% mdparagraph:P008:begin

\textbf{B1. cuman, “come”.} \textbf{C} is \textbf{k}, as in \emph{kit}. It is a stopped sound, without the \textbf{ch} beginning heard in \emph{chin}. The English gloss illustrates the initial consonant, not the whole Old English word.

% mdparagraph:P008:end



% answer: PRON-P02b

% mdparagraph:P009:begin

\textbf{B2. cyrice, “church”.} Both \textbf{c} letters have the \textbf{ch} sound of \emph{chin}. Keep three syllables and the light final \textbf{e}. Short \textbf{y} remains the rounded-front vowel taught in the guide; the modern English gloss is not a complete pronunciation model.

% mdparagraph:P009:end



% answer: PRON-P02c

% mdparagraph:P010:begin

\textbf{B3. cēlan, “cool”.} \textbf{C} is \textbf{k}. This is a supplied exception to the front-vowel tendency. The following \textbf{ē} does not authorize changing it to \textbf{ch}.

% mdparagraph:P010:end



% answer: PRON-P02d

% mdparagraph:P011:begin

\textbf{B4. giefan, “give”.} Initial \textbf{g} is consonant \textbf{y}, as in \emph{yes}. The following \textbf{ie} is the taught single-syllable glide beginning with the \emph{see} vowel. Neither hard \textbf{g} nor the sound of English \textbf{j} is the requested initial sound.

% mdparagraph:P011:end



% answer: PRON-P02e

% mdparagraph:P012:begin

\textbf{B5. fugol, “bird”.} \textbf{G} is continuous and voiced: bring the back of the tongue near the hard-\textbf{g} position, leave a gap, and sustain the airflow with throat vibration. A complete closure would give a hard stop; omitting the sound would lose a consonant.

% mdparagraph:P012:end



% answer: PRON-P02f

% mdparagraph:P013:begin

\textbf{B6. scip, “ship”.} \textbf{Sc} is the single \textbf{sh} sound of \emph{ship}. Do not pronounce separate \textbf{s} and \textbf{k} here. That separate pair belongs to the supplied exception \textbf{frosc}, “frog”.

% mdparagraph:P013:end



% answer: PRON-P02g

% mdparagraph:P014:begin

\textbf{B7. niht, “night”.} Use the front, unvoiced continuous \textbf{h} sound: raise the middle of the tongue toward the roof of the mouth and let breath pass through the gap. Both \textbf{h} and \textbf{t} are audible. The modern gloss does not license a silent \textbf{h} or a long vowel.

% mdparagraph:P014:end



% answer: PRON-P02h

% mdparagraph:P015:begin

\textbf{B8. wyrt, “herb; plant”.} \textbf{Y} is a short rounded-front vowel. Keep the tongue position of \emph{bit} while rounding the lips as for \textbf{oo} in \emph{food}. Do not pull the tongue back into an ordinary \textbf{oo} sound or make the consonant \textbf{y} of \emph{yes}.

% mdparagraph:P015:end



\reserve{5\baselineskip}

% mdparagraph:P016:begin

\section*{C. Voice}

% mdparagraph:P016:end



% answer: PRON-P03a

% mdparagraph:P017:begin

\textbf{C1. cnīf, “knife”.} Final \textbf{f} is unvoiced, as in \emph{fan}. The preceding vowel does not voice a word-final \textbf{f} under this guide's rule. The initial \textbf{c} is also sounded, as \textbf{k}.

% mdparagraph:P017:end



% answer: PRON-P03b

% mdparagraph:P018:begin

\textbf{C2. cnīfas, “knives”.} \textbf{F} is voiced, as in \emph{van}, between \textbf{ī} and \textbf{a}. This is the same written letter as in \textbf{cnīf}, now in a different sound environment. The final \textbf{s} remains unvoiced.

% mdparagraph:P018:end



% answer: PRON-P03c

% mdparagraph:P019:begin

\textbf{C3. rīsan, “rise”.} \textbf{S} is voiced, as in \emph{zip}, because it stands between voiced vowels. Its spelling does not require the unvoiced sound of \emph{sip} in this position.

% mdparagraph:P019:end



% answer: PRON-P03d

% mdparagraph:P020:begin

\textbf{C4. hūs, “house”.} Final \textbf{s} is unvoiced, as in \emph{sip}. It remains unvoiced despite the preceding vowel. A voiced \textbf{z} is not the selected word-final reading.

% mdparagraph:P020:end



% answer: PRON-P03e

% mdparagraph:P021:begin

\textbf{C5. ōðer, “other”.} \textbf{Ð} is voiced, as in \emph{this}, between vowels. The sound follows its environment here, not a distinction between the letters \textbf{ð} and \textbf{þ}.

% mdparagraph:P021:end



% Keep the complete exception explanation together.
\reserve{6\baselineskip}

% answer: PRON-P03f

% mdparagraph:P022:begin

\textbf{C6. hǣlðu, “health”.} \textbf{Ð} is unvoiced, as in \emph{thin}. Although \textbf{l} and \textbf{u} carry voice, this word is an explicitly supplied exception. Applying the ordinary surrounding-sounds rule without its stated limits would give the wrong target.

% mdparagraph:P022:end



\reserve{5\baselineskip}

% mdparagraph:P023:begin

\section*{D. Syllables and stress}

% mdparagraph:P023:end



% answer: PRON-P04a

% mdparagraph:P024:begin

\textbf{D1. habban, “have”.} Two syllables; stress the first. \textbf{Bb} is prolonged: hold the lip closure before releasing into the second syllable. There is no extra vowel or third syllable between the two written consonants.

% mdparagraph:P024:end



% answer: PRON-P04b

% mdparagraph:P025:begin

\textbf{D2. fugol, “bird”.} Two syllables; stress the first. There is no doubled consonant. Keep the second vowel audible and the medial \textbf{g} continuous and voiced; dropping the unstressed \textbf{o} is not the requested careful reading.

% mdparagraph:P025:end



% answer: PRON-P04c

% mdparagraph:P026:begin

\textbf{D3. heofonan, “heaven”.} Three syllables; stress the first. There is no doubled consonant. \textbf{Eo} forms one glide, followed by syllables containing \textbf{o} and \textbf{a}. Keep both unstressed vowels audible; copying the two syllables of modern \emph{heaven} would lose a syllable.

% mdparagraph:P026:end



% answer: PRON-P04d

% mdparagraph:P027:begin

\textbf{D4. gescēop, “created”.} Two syllables; stress the second, \textbf{scēop}. There is no doubled consonant. The prefix \textbf{ge-} is unstressed; long \textbf{ēo} remains one vowel glide within the stressed syllable. Two vowel letters do not create two syllables there.

% mdparagraph:P027:end



\reserve{5\baselineskip}

% mdparagraph:P028:begin

\section*{E. A sentence adapted from Cook}

% mdparagraph:P028:end



% mdparagraph:P029:begin

The supplied meaning is “God created heaven and earth.” No independent translation is tested.

% mdparagraph:P029:end



% answer: PRON-P05a

% mdparagraph:P030:begin

\textbf{E1. gescēop, “created”.} \textbf{G} has the sound of \textbf{y} in \emph{yes}; \textbf{sc} has \textbf{sh}. The prefix contains a light \textbf{e}; the stressed second syllable contains long \textbf{ēo}, moving from a held, steady \emph{day}-type vowel toward \textbf{oo}. Keep the prefix vowel: \textbf{g} and \textbf{sc} are not adjacent sounds with no vowel between them.

% mdparagraph:P030:end



% answer: PRON-P05b

% mdparagraph:P031:begin

\textbf{E2. heofonan, “heaven”.} \textbf{F} is voiced, as in \emph{van}, between the \textbf{eo} glide and the following \textbf{o}, both voiced vowels. Keep three syllables, with stress on the first. The rule depends on surrounding sounds, not on the modern English spelling of the gloss.

% mdparagraph:P031:end



% answer: PRON-P05c

% mdparagraph:P032:begin

\textbf{E3. eorðan, “earth”.} \textbf{Ð} is voiced, as in \emph{this}, between voiced \textbf{r} and \textbf{a}. The rule is “between voiced sounds”, not only “between vowels”. Keep the initial \textbf{eo} glide and the second syllable; the single-syllable English gloss is not a pronunciation model.

% mdparagraph:P032:end



\end{document}
